\section{探索与利用是同一笔预算}
\label{sec:17-Constrained-Guarantees/04-Online-Learning-and-Regret/03}

一个落地页准备了六个版本，每来一位访客只能展示其中一个。展示了版本 A，就永远不会知道这位访客看到版本 B 会不会下单——那个人已经走了，不会再来一次。

这和上一节的八个信号源差在哪里？信号源每天都会亮出各自的对错，哪怕当天没跟它走，第二天照样能把它的损失记进 \(L_{i,t}\)。落地页不会。第 \(t\) 轮结束后，我们只拿到 \(\ell_t(i_t)\) 这一个数，其余五个版本那一轮的表现是空白。

Hedge 的更新式在这里直接卡住：\(p_{i,t}\propto e^{-\eta L_{i,t-1}}\) 要求每一轮都能给全部 \(N\) 个候选各记一笔账，而现在只有一笔可记。在线梯度下降同样卡住，它要的是当前点上的次梯度，可我们连当前决策以外的函数值都看不到。

这一处改动叫作 bandit 反馈，得名于老虎机——扳下哪个拉杆，就只看到哪个拉杆吐出多少。改的只是信息通道，问题的其余部分一个字没动。而最优遗憾的量级立刻从 \(\sqrt{T\ln N}\) 变成 \(\sqrt{TN}\)。

\subsection{\texorpdfstring{\(\ln N\) 变成 \(N\)}{ln N 变成 N}}

把两个量级并排：全信息下 \(\Theta(\sqrt{T\ln N})\)，bandit 下 \(\Theta(\sqrt{TN})\)。上界由 EXP3 一类算法给出（严格写是 \(O(\sqrt{TN\ln N})\)），下界是 \(\Omega(\sqrt{TN})\)，两者只差一个对数因子。

对数变成一次方，这个差别在数值上很直白。候选从 \(10\) 个增到 \(100\) 个，全信息下的界放大 \(\sqrt{\ln100/\ln10}\approx1.41\) 倍，bandit 下放大 \(\sqrt{100/10}\approx3.16\) 倍。候选从 \(10\) 增到 \(10^{6}\)，前者放大 \(2.4\) 倍，后者放大 \(316\) 倍。全信息里``多准备几个候选''几乎免费，bandit 里它是实打实的开销。

这个差别的来源可以指出来。EXP3 仍然想用指数权重，可它记不下 \(\ell_t(i)\)，只好造一个替代品：

\[
\hat\ell_t(i)=\frac{\ell_t(i)\,\ind\{i_t=i\}}{p_{i,t}}.
\]

只有被选中的那个候选有非零值，除以选中它的概率是为了补偿。这个估计是无偏的——\(\E[\hat\ell_t(i)]=\ell_t(i)\)，期望对 \(i_t\sim p_t\) 取。但它的二阶矩是 \(\ell_t(i)^2/p_{i,t}\)，概率越小方差越大。把 Hedge 证明里那一步重做，\(\eta^2/8\) 那个位置会换成一个与 \(\sum_i p_{i,t}\E[\hat\ell_t(i)^2]\) 有关的量，而它是 \(N\) 的量级而不是常数。\(N\) 就是从这里进来的。

重要性权重的这套记账在别处出现过：\hyperref[sec:16-Control-and-Reinforcement-Learning/02-Reinforcement-Learning/08]{``Off-policy 校正用权重交换分布偏差与方差''}里，用一套策略采的数据去评估另一套策略，靠的是同一个除以概率的动作，付出的也是同一种方差膨胀；\hyperref[sec:05-Probability/08-Limit-Theorems/05]{``Monte Carlo：从随机样本到数值答案''}里的重要性采样也是这个式子，那里同样是提议分布在某处过小就会让估计不可用。信息缺失在这里被量化成了一个具体的数：方差里那个 \(1/p\)。

由此也能看出探索为什么是强制的。若某个候选的概率 \(p_{i,t}\) 被压到接近零，它的损失估计方差就会炸开，界随之失效。算法必须持续给每个候选留一份不小的概率，哪怕当前的证据说它不好。这份留出来的概率就是探索的预算，它与利用共用同一个概率单纯形——给探索多一分，给当前最优就少一分。两者不是两笔账，是同一笔。

\subsection{乐观地对待还没看清的选项}

上面的设定里对手是恶意的。若换成一个温和些的假设——每个候选 \(i\) 的奖励独立同分布地来自一个均值为 \(\mu_i\) 的固定分布——就能做得好得多。这是随机 bandit，UCB 是它的标准解法。

想法是给每个候选的经验均值加一条置信半径：

\[
\mathrm{UCB}_i(t)=\hat\mu_i(t)+\sqrt{\frac{2\ln t}{n_i(t)}},
\]

其中 \(n_i(t)\) 是候选 \(i\) 到目前为止被选中的次数，然后贪心地选 \(\mathrm{UCB}\) 最大的那个。半径的形状来自 Hoeffding 界（\hyperref[sec:05-Probability/07-Transforms-and-Bounds/04]{``Chernoff 方法与 Hoeffding 集中界''}）：\(n\) 个有界样本的均值偏离真值超过 \(\sqrt{2\ln t/n}\) 的概率是 \(t^{-4}\) 量级，再对所有候选与所有轮次用一次联合界（\hyperref[sec:05-Probability/07-Transforms-and-Bounds/02]{``从 Union Bound 到 Markov 与 Chebyshev''}），全程失效的概率仍然可控。

这条规则同时处理了两件事，而只用了一个式子。均值高的候选 \(\mathrm{UCB}\) 高，这是利用；被试得少的候选半径大、\(\mathrm{UCB}\) 也高，这是探索。一个候选只有在``均值确实低''且``已经试够多次''时才会被冷落。这句话常被写成``乐观面对不确定性''：不确定的时候按最好的可能情况估值，于是要么它真的好、我们赚到，要么它其实不好、我们把它试清楚了。两种结局都不亏。

遗憾界的形式是

\[
R_T=O\Bigl(\sum_{i:\Delta_i>0}\frac{\ln T}{\Delta_i}\Bigr),
\qquad \Delta_i=\mu^\star-\mu_i .
\]

\(\Delta_i\) 是候选 \(i\) 与最优候选的差距，它出现在分母上。推导这个位置只需要一行：要把候选 \(i\) 从最优候选中区分开，置信半径得缩到 \(\Delta_i\) 以下，也就是需要 \(n_i\approx\ln T/\Delta_i^2\) 次尝试；每一次尝试比选最优多损失 \(\Delta_i\)；两者相乘就是 \(\ln T/\Delta_i\)。

于是``越难区分的候选，越贵''。两个转化率相差 \(5\) 个百分点的版本，试几十次就能分开；相差 \(0.5\) 个百分点的版本，要试上百倍的次数才分得开，而每次试错本身也在扣分。这与随机对照实验里样本量随效应量平方反比增长是同一条算术。

分母上的 \(\Delta_i\) 还解释了 \(\ln T\) 与 \(\sqrt T\) 怎样相处。间隙固定时遗憾只按 \(\ln T\) 增长，远好过 \(\sqrt T\)；但把间隙也交给对手挑，\(\Delta\) 会被调到最难受的位置：总遗憾大致是 \(\min\{T\Delta,\;N\ln T/\Delta\}\)，两支在 \(\Delta\approx\sqrt{N\ln T/T}\) 处相交，取值 \(\sqrt{TN\ln T}\)。间隙相关的界与最坏情况的界在这里握手，并不矛盾——前者是给定 \(\Delta\) 之后的话，后者是对 \(\Delta\) 取最坏之后的话。

\begin{center}
\begin{tikzpicture}[x=1cm,y=1cm,font=\small,>=stealth]
  \draw[->,black!60] (-0.15,0) -- (7.1,0) node[right,font=\footnotesize] {\(T\)};
  \draw[->,black!60] (0,-0.15) -- (0,4.9) node[above,font=\footnotesize] {\(R_T\)};
  \foreach \x/\lab in {0/{10^1},1.3/{10^2},2.6/{10^3},3.9/{10^4},5.2/{10^5},6.5/{10^6}}
    \draw[black!45] (\x,0) -- (\x,-0.12) node[below,font=\scriptsize] {\(\lab\)};
  \foreach \y/\lab in {1/{10^1},2/{10^2},3/{10^3},4/{10^4}}
    \draw[black!45] (0,\y) -- (-0.12,\y) node[left,font=\scriptsize] {\(\lab\)};
  \draw[black!15] (0,1) -- (6.5,1);
  \draw[black!15] (0,2) -- (6.5,2);
  \draw[black!15] (0,3) -- (6.5,3);
  \draw[black!15] (0,4) -- (6.5,4);
  \draw[blue!65!black,thick] (0,0.68) -- (6.5,3.18);
  \draw[red!70!black,thick] (0,1.65) -- (6.5,4.15);
  \draw[green!45!black,thick]
    (0,1.57) -- (1.3,1.87) -- (2.6,2.04) -- (3.9,2.17) -- (5.2,2.27) -- (6.5,2.34);
  \draw[green!45!black,thick,dashed]
    (0,2.27) -- (1.3,2.57) -- (2.6,2.74) -- (3.9,2.87) -- (5.2,2.97) -- (6.5,3.04);
  \draw[black!45] (6.5,3.18) -- (6.62,3.42);
  \draw[black!45] (6.5,3.04) -- (6.62,2.80);
  \node[anchor=west,font=\footnotesize,red!70!black] at (6.66,4.15) {bandit \(\sqrt{TN}\)};
  \node[anchor=west,font=\footnotesize,blue!65!black] at (6.66,3.42) {全信息 \(\sqrt{T\ln N}\)};
  \node[anchor=west,font=\footnotesize,green!45!black] at (6.66,2.80) {UCB，间隙小};
  \node[anchor=west,font=\footnotesize,green!45!black] at (6.66,2.30) {UCB，间隙大};
\end{tikzpicture}

\emph{两轴都是对数刻度。看得见的只有两件事：少一条信息通道，整条线往上平移一个量级（蓝到红）；间隙缩小十倍，那条近乎水平的线也往上平移一个量级（实绿到虚绿）。信息越少越贵，区分越难越贵。}
\end{center}

图里绿色那两条几乎是平的，因为 \(\ln T\) 在对数坐标下几乎不动。间隙小的那一条在早期还高过红线——这不是画错了：间隙相关的界在轮次不够多时确实弱于最坏情况的界，此时后者才是有效的那一个。两条界同时成立，取小的用。

条件对账。UCB 的置信半径是照着有界或次高斯奖励算的；奖励重尾时那个 \(\sqrt{2\ln t/n}\) 严重低估了偏差，算法会长期锁定在一个被几次幸运抽样抬高了均值的候选上，遗憾退化为线性。修补要么截断奖励，要么换成基于中位数的估计。第二个条件是同一候选的多次奖励之间独立同分布；若展示这个版本会改变访客后续的行为，问题就已经不是 bandit 而是 MDP，动作影响状态、状态影响未来收益，得换成\hyperref[sec:16-Control-and-Reinforcement-Learning/02-Reinforcement-Learning/01]{``MDP 把决策问题写成随机过程''}那套框架。第三个条件容易被忽略：\(\hat\mu_i\) 是在自适应采样下算出来的，被选中的次数本身依赖历史，所以它不是一列独立样本的简单平均，严格的证明要用鞅形式的集中不等式而不是直接套 Hoeffding。第四个条件是环境不对抗；把随机 bandit 换成对抗 bandit，UCB 的遗憾可以是线性的，那里必须用 EXP3 那种带强制探索的随机化算法。

还有一条与本书别处呼应的边界。遗憾衡量的是累计收益，它不保证最终能\emph{指认}出哪个候选最好。想要``以 \(95\%\) 的把握说出最优版本''是另一个目标，叫最优臂识别，最优算法与 UCB 不同。把这两件事混起来，是在线实验里常见的口径错误——\hyperref[sec:13-Machine-Learning/01-Learning-Problems-and-Evaluation/02]{``损失、风险、目标与指标不是同一个量''}讲的正是这类混淆，而\hyperref[sec:09-Convex-Optimization/08-Nonconvex-and-Stochastic-Optimization/06]{``优化误差、统计误差与泛化不能混为一谈''}提醒的是同一种分账习惯。同理，用 bandit 分配流量之后，日志里各版本的样本量不再是随机分配的结果，直接拿它做显著性检验会得到错误的 \(p\) 值。

\subsection{当对手也在学习}

到这里为止，对手一直被当成一个不需要解释的东西：它可以任意出招，我们只保证不比事后最优的固定决策差。这个设定足够安全，也足够悲观。

真实场景里，对面往往不是一个无所不知的恶意实体，而是另一个和我们处境相同的学习者。反欺诈那边有人在跑他自己的优化循环，广告竞价的另一侧是别家的出价算法，定价系统对面是同样在调价的竞争对手。双方都在读对方的行为，都在更新自己的策略。

一旦两边都用无遗憾算法，前面的分析就开始变形。遗憾的定义里那个``事后最优固定决策''，是相对于对手实际打出的那串损失函数算的；而对手的出招取决于我们的行为，我们的行为又取决于对手先前的出招。两边的基准都在动，谁也不是谁的固定背景。

这时该问的不再是``我们收敛到哪个最优点''——最优点这个概念已经没有确定的所指了。可以问的是：这套互相追逐的动力学长期会稳定在什么状态？如果它稳定下来，稳定在什么位置？如果不稳定，不稳定成什么样子？

答案不是一个点，是一种平衡：双方各自都不再有单方面改进的余地。这是下一个单元的对象。
